\documentclass[twocolumn]{aastex631}
\begin{document}
\title{A residual reionization ionizing-photon-budget shortfall for star-forming galaxies at z\~{}9}
\author{NebulaMind Lab (autonomous pipeline)}
\affiliation{NebulaMind Lab, \url{https://lab.nebulamind.net}}
\begin{abstract}
Reionization ionizing-photon-budget reconciliation at z\~{}9: star-forming galaxies require f\_esc=0.452 (+0.455/-0.231) to close the budget (Madau-Dickinson SFRD, log xi\_ion=25.5+/-0.15, clumping C=2-5, JWST-SFRD tail), versus indirect-proxy-inferred f\_esc=0.062 (+0.108/-0.039) from LzLCS O32/beta calibrations. Median delta(required-inferred)=+0.362 dex-frac (16-84\%: +0.117 to +0.818); 94\% of the systematic MC shows a shortfall. Conclusion: a genuine SHORTFALL remains, and the sign holds under both O32 and beta calibrations. The result is bounded by the xi\_ion x clumping x proxy-calibration systematic, not by statistics. Generated autonomously from public data (jwst) via the NebulaMind Lab runner: a bounded, reproducible, descriptive study.
\end{abstract}
\section{Introduction}
Introduction: Recent studies have highlighted a potential crisis in the photon budget for reionization, suggesting that current estimates of ionizing photons from star-forming galaxies may not be sufficient to account for the observed reionization history [Muoz2024]. This discrepancy has led to increased scrutiny of the assumptions and parameters used in these calculations, such as the escape fraction (f\_esc) of ionizing photons from galaxies. Previous work has shown that the galaxy ionizing photon budget is highly sensitive to the choice of f\_esc and other factors, including the clumping factor (C) and the ionization efficiency (xi\_ion) [Davies2021]. To address this issue, we revisit the reionization-photon-budget calculation using a literature-anchored approach.
\section{Data and method}
Data and Method: We adopt the cosmic star formation rate density (SFRD) from Madau \& Dickinson (2014), which provides an analytic fitting function for the SFRD based on a compilation of observational data. For the ionizing photon production efficiency (xi\_ion), we use published values with log xi\_ion = 25.5 ± 0.15 [Madau2017]. The escape fraction is inferred using the O32/beta proxy calibration from LzLCS, which provides a relationship between the optical emission line ratio and f\_esc [Chisholm+22, Flury+22; Simmonds+24]. We also explore the impact of varying clumping factors (C=2-5) on our results.
\section{Result}
Result: Our calculations show that star-forming galaxies require an escape fraction of f\_esc = 0.452 (+0.455/-0.231) to close the reionization-photon-budget at z\~{}9, assuming the Madau-Dickinson SFRD and log xi\_ion=25.5±0.15. However, indirect-proxy-inferred values from LzLCS O32/beta calibrations yield a significantly lower f\_esc = 0.062 (+0.108/-0.039). The median difference between the required and inferred escape fractions is +0.362 dex-frac (16-84\%: +0.117 to +0.818), with 94\% of our systematic Monte Carlo simulations indicating a shortfall in ionizing photons.
\begin{figure}\centering\includegraphics[width=\columnwidth]{result.png}\caption{A residual reionization ionizing-photon-budget shortfall for star-forming galaxies at z\~{}9}\end{figure}
\section{Caveats}
Caveats: Our analysis relies on several assumptions and simplifications, which may introduce uncertainties into our results. For example, we assume a fixed clumping factor (C=2-5) and do not account for potential variations in this parameter across different environments or redshifts. Additionally, our use of the O32/beta proxy calibration is based on a specific set of observational data and may not capture the full range of galaxy properties. Furthermore, our calculations are sensitive to the choice of xi\_ion and SFRD, which are subject to their own uncertainties and biases. Finally, our analysis does not incorporate any new observational data or account for potential systematic errors in existing datasets.
\section*{References}
\footnotesize
[Muoz2024] Muñoz et al. (2024). Reionization after JWST: a photon budget crisis?. bibcode 2024MNRAS.535L..37M \par [Davies2021] Davies et al. (2021). The Predicament of Absorption-dominated Reionization: Increased Demands on Ionizing Sources. bibcode 2021ApJ...918L..35D \par [Park2022] Park et al. (2022). Calibrating excursion set reionization models to approximately conserve ionizing photons. bibcode 2022MNRAS.517..192P \par [Duncan2015] Duncan et al. (2015). Powering reionization: assessing the galaxy ionizing photon budget at z < 10. bibcode 2015MNRAS.451.2030D \par [Madau2017] Madau (2017). Cosmic Reionization after Planck and before JWST: An Analytic Approach. bibcode 2017ApJ...851...50M

\end{document}
